\section{State of Research}

The piezoresistive effect has been extensively investigated in both experimental and engineering contexts, 
particularly in the development of semiconductor-based sensors and microelectromechanical systems (MEMS) \cite{smith1954,doll2013}. 
Early studies established the strong dependence of electrical conductivity on mechanical strain in semiconductors such as silicon, 
highlighting the role of crystal structure and carrier transport in determining the sensitivity of the effect.

From a materials physics perspective, piezoresistivity originates from strain-induced modifications of the electronic structure 
and charge transport properties. In semiconductors, mechanical deformation can alter band structure, carrier mobility, and scattering mechanisms, 
leading to significant changes in conductivity. These effects have been studied using both phenomenological models and more detailed 
electronic transport theories.

In practical applications, however, the coupling between mechanical deformation and electrical response is often modeled in a simplified 
and application-specific manner, where the dependence of conductivity on strain is introduced empirically. While such approaches are effective 
for device design and optimization, they do not always provide a transparent connection between the underlying physical mechanisms and the 
resulting macroscopic behavior.

At the continuum level, models based on partial differential equations are widely used to describe both elasticity and electrical conduction. 
However, the integration of strain-dependent transport properties into such models is typically treated in a phenomenological way, and a 
systematic computational framework that bridges physical insight and continuum modeling remains an open challenge.

Such a framework is particularly relevant for understanding emerging materials and 
devices where electromechanical coupling plays a central role.

Based on these considerations, the proposed study is guided by the following research questions:

\begin{itemize}
    \item How can strain-dependent conductivity be incorporated into a consistent continuum model of electromechanical coupling?
    \item How does mechanical deformation influence the spatial distribution of electrical potential in conductive materials?
    \item What qualitative and quantitative effects arise from the nonlinear coupling between mechanics and electrical transport?
    \item To what extent can flexible (e.g. data-driven or physics-informed) approaches improve the modeling of $\sigma(\varepsilon)$?
\end{itemize}

To address these questions, the research will be structured around the following key components:

\begin{itemize}
    \item Physical background and modeling assumptions
    \item Formulation of the coupled electromechanical model
    \item Analysis of the nonlinear system
    \item Computational methods and numerical experiments
    \item Interpretation of results in the context of material behavior
\end{itemize}

